\documentclass{article}
\usepackage[preprint]{neurips_2026}
\usepackage{amsmath,amssymb,graphicx,booktabs,hyperref,xcolor,algorithm,algpseudocode,subcaption,CJK,pifont}
\newcommand{\cmark}{\ding{51}}
\newcommand{\xmark}{\ding{55}}

\title{k\=odo: Persistent Structured Memory with Real-Time Cross-Session Knowledge Sharing for LLM Coding Agents}
\author{Xuan J.}
\date{}

\begin{document}
\maketitle

\begin{abstract}
Large Language Model (LLM) coding agents suffer from complete session amnesia---every conversation starts from a blank slate, causing repeated mistakes, forgotten conventions, and redundant re-discovery of project knowledge. We present \textbf{k\=odo}, a lightweight persistent memory layer that implements a \emph{memory taxonomy} of six typed categories (conventions, mistakes, decisions, preferences, patterns, notes) stored in a local SQLite database with FTS5 full-text search, exposed to agents via the Model Context Protocol (MCP), with real-time cross-session knowledge sharing through a Unix domain socket pub/sub hub. We further introduce \emph{self-evolving memory}---an automated prune/merge/promote cycle that drives \emph{memory evolution} to maintain recall precision as the store grows. In controlled Monte Carlo simulations across 30 independent trials with 50 sessions each, k\=odo reduces cross-session mistake repetition by 27.4\% ($p < 10^{-6}$), improves convention adherence by 53.3\% ($p < 10^{-6}$), and achieves sub-millisecond cross-terminal knowledge transfer ($0.31 \pm 0.01$\,ms vs.\ $183 \pm 9.8$\,s manual re-discovery, a $5.8 \times 10^5\times$ speedup). Ablation studies across six configurations confirm that typed structured memory with FTS5 achieves Precision@5 of $0.816 \pm 0.011$, outperforming unstructured logs ($0.413 \pm 0.018$, $+97.6\%$) and vector embeddings ($0.577 \pm 0.011$, $+41.4\%$) for coding-specific retrieval. Self-evolving memory further improves precision to $0.878 \pm 0.006$ ($+7.6\%$) while reducing store size by 60\% over 10 evolution cycles. The system requires approximately 400 lines of JavaScript with zero cloud dependencies.
\end{abstract}

\section{Introduction}

AI coding agents powered by large language models are increasingly used for software development, from generating boilerplate to debugging complex systems~\cite{chen2021codex,li2023starcoder,roziere2023codellama,jimenez2024swebench,yang2024sweagent}. Yet they suffer from a fundamental limitation: \emph{session amnesia}. Every new conversation starts with zero context about the project's conventions, past mistakes, or architectural decisions. A developer who corrects an agent's import style in one session will encounter the same error in the next. A bug fix discovered in terminal A remains invisible to the agent in terminal B.

Consider a concrete scenario: a developer uses an AI agent to implement a REST API. In session~1, the agent generates Express handlers using CommonJS \texttt{require()} syntax; the developer corrects it to ESM \texttt{import}. In session~2, the agent generates middleware---again with \texttt{require()}. The developer corrects it again. In session~3, working in a second terminal on the same project, a colleague's agent generates test utilities---once more with \texttt{require()}. Each correction takes 2--5 minutes, and across a week of development, these redundant corrections accumulate into hours of wasted effort. This is not a hypothetical: it is the daily experience of developers using stateless AI coding agents.

This problem is not merely an inconvenience---it represents a systematic failure mode. We identify three concrete manifestations:

\begin{enumerate}
    \item \textbf{Mistake repetition}: Without memory of prior corrections, agents repeat the same errors across sessions, wasting developer time on redundant fixes.
    \item \textbf{Convention drift}: Project-specific coding standards (import styles, error handling patterns, naming conventions) must be re-taught every session, leading to inconsistent codebases.
    \item \textbf{Knowledge silos}: In multi-terminal workflows, discoveries in one session are invisible to others, requiring manual re-discovery or copy-paste of context.
\end{enumerate}

Existing approaches to agent memory fall short for coding-specific use cases. Generic memory systems like mem0~\cite{mem0} store untyped text blobs without coding-domain structure. Agent-specific solutions like claude-mem~\cite{claudemem} are locked to a single agent. RAG-based approaches using vector embeddings~\cite{lewis2020rag} optimize for semantic similarity but lack the structured type system needed to distinguish a ``convention'' from a ``mistake'' from a ``decision.'' MemGPT~\cite{packer2023memgpt} introduces virtual memory paging for long conversations but does not persist across sessions. None of these systems provide real-time cross-session knowledge sharing or self-evolving memory management.

We present \textbf{k\=odo} (\begin{CJK}{UTF8}{min}コード\end{CJK}, Japanese for ``code''), a persistent structured memory layer designed specifically for LLM coding agents. k\=odo introduces a \emph{memory taxonomy}---a principled categorization of coding knowledge into six functional types---and implements \emph{memory evolution} through automated lifecycle management. Our three key contributions are:

\begin{enumerate}
    \item \textbf{Typed memory with full-text search}: A SQLite-backed store with six memory types (convention, mistake, decision, preference, pattern, note) and FTS5 indexing, enabling precise retrieval that outperforms both unstructured logs and vector embeddings (\S\ref{sec:ablation}).
    \item \textbf{Cross-session knowledge sharing}: A Unix domain socket pub/sub hub that broadcasts memories between parallel agent sessions in sub-millisecond time, plus a static export system for agent-native configuration files (\S\ref{sec:transfer}).
    \item \textbf{Self-evolving memory}: An automated prune/merge/promote cycle based on access patterns that improves recall precision over time while compacting the store (\S\ref{sec:evolve}).
\end{enumerate}

k\=odo is agent-agnostic, operating as a standard MCP server compatible with Claude Code, Cursor, Kiro, Codex CLI, and any MCP-compliant client. The entire system is approximately 400 lines of JavaScript with zero cloud dependencies.

\section{Related Work}

\paragraph{Agent memory systems.} MemGPT~\cite{packer2023memgpt} introduces a virtual memory hierarchy for LLM agents, paging context in and out of the limited context window. While effective for long conversations, it does not persist across sessions or share between agents. mem0~\cite{mem0} provides a generic memory layer with embedding-based retrieval, but stores untyped text without coding-domain structure. claude-mem~\cite{claudemem} captures session transcripts for Claude Code but is locked to a single agent and lacks structured retrieval. MemoryBank~\cite{zhong2024memorybank} enhances LLMs with long-term memory using Ebbinghaus forgetting curves, but relies on embedding similarity without domain-specific typing. Hu et~al.~\cite{hu2024memoryaugmented} survey the broader landscape of memory-augmented LLMs, identifying persistent external stores as a key open challenge. None of these systems implement a \emph{memory taxonomy} that distinguishes the functional role of different knowledge types.

\paragraph{Retrieval-augmented generation.} RAG systems~\cite{lewis2020rag,borgeaud2022retro} retrieve relevant documents to augment LLM context. While powerful for general knowledge, they rely on embedding similarity which conflates semantically similar but functionally distinct memories (e.g., a convention about imports vs.\ a mistake about imports). Gao et~al.~\cite{gao2023retrieval} survey the RAG landscape and identify structured metadata as a key complement to dense retrieval---a finding that motivates k\=odo's typed memory design. k\=odo's type system provides orthogonal filtering that complements text-based retrieval.

\paragraph{Tool-augmented agents.} The Model Context Protocol (MCP)~\cite{mcp2024} standardizes how applications provide context to LLMs through tool interfaces. k\=odo leverages MCP to expose memory operations (remember, recall, forget, stats) as first-class tools that agents can invoke during sessions, enabling a closed-loop learning cycle. This builds on the broader trend of tool-augmented language models~\cite{schick2023toolformer,yao2023react,zhang2024codeagent}, where external tools extend the capabilities of LLMs beyond their parametric knowledge. The attention mechanism~\cite{vaswani2017attention} underlying modern LLMs~\cite{brown2020gpt3,touvron2023llama} provides powerful in-context learning, but this learning is ephemeral---it vanishes when the context window resets.

\paragraph{Generative agents and persistent memory.} Park et~al.~\cite{park2023generative} demonstrate that generative agents with memory retrieval, reflection, and planning can exhibit believable human behavior in simulated environments. Their memory stream architecture inspires k\=odo's approach of storing timestamped observations with relevance-based retrieval, though we specialize the design for coding-domain knowledge with typed memories and full-text search rather than embedding-based retrieval. Reflexion~\cite{shinn2023reflexion} enables agents to learn from verbal self-reflection, and Voyager~\cite{wang2024voyager} maintains a skill library for open-ended exploration---both demonstrating the value of persistent agent memory, though neither addresses cross-session sharing or coding-specific type systems. Multi-agent frameworks like ChatDev~\cite{qian2024chatdev} and MetaGPT~\cite{hong2024metagpt} coordinate multiple LLM agents for software development but rely on in-context communication rather than persistent shared memory, limiting knowledge transfer to a single execution.

\paragraph{Knowledge management in software engineering.} Prior work on organizational memory in software teams~\cite{robillard2004effective,ko2007information} identifies knowledge loss during developer transitions as a major cost. SWE-bench~\cite{jimenez2024swebench} evaluates agents on isolated GitHub issues but does not measure cross-session learning. The bootstrap methodology~\cite{efron1993bootstrap} and Welch's $t$-test~\cite{welch1947} provide the statistical foundations for our experimental evaluation. k\=odo addresses the analogous problem for AI agents: institutional knowledge that should persist across sessions and transfer between tools.

\paragraph{Code generation models.} The rapid progress in code-specialized LLMs---from Codex~\cite{chen2021codex} to StarCoder~\cite{li2023starcoder} and Code Llama~\cite{roziere2023codellama}---has made AI coding agents increasingly capable at individual tasks. However, these models operate statelessly: each invocation is independent, with no mechanism to accumulate project-specific knowledge over time. k\=odo provides the missing persistence layer that enables these models to learn from their deployment history. SQLite's FTS5 extension~\cite{sqlite2024} provides the efficient full-text indexing that makes this practical at scale.

\section{Method}
\label{sec:method}

\subsection{Architecture Overview}

k\=odo consists of four components: (1) a typed memory store backed by SQLite with FTS5, (2) an MCP server exposing memory operations to agents, (3) a cross-terminal pub/sub hub for real-time knowledge sharing, and (4) a self-evolving memory manager that implements memory evolution through automated lifecycle management. The design philosophy prioritizes simplicity and zero-configuration deployment: the entire system is a single Node.js process with no external dependencies beyond SQLite. Figure~\ref{fig:architecture} illustrates the architecture.

\begin{figure}[t]
    \centering
    \includegraphics[width=0.85\textwidth]{figures/architecture.pdf}
    \caption{k\=odo architecture. Agents interact with the MCP server, which reads/writes a typed SQLite store with FTS5 indexing. The cross-terminal hub broadcasts new memories to all connected sessions via Unix domain sockets. The evolve daemon periodically prunes, merges, and promotes memories.}
    \label{fig:architecture}
\end{figure}

\subsection{Typed Memory Store}

Each memory $m$ is a tuple $(t, c, \tau, s, p, a, r)$ where $t \in \mathcal{T}$ is the memory type, $c$ is the content string, $\tau$ is a set of tags, $s$ is the source (manual or agent), $p$ is the project identifier, $a \in \mathbb{N}$ is the access count, and $r \in \mathbb{R}^+$ is the relevance score. The type system $\mathcal{T}$ consists of six categories:

\begin{equation}
    \mathcal{T} = \{\texttt{convention}, \texttt{mistake}, \texttt{decision}, \texttt{preference}, \texttt{pattern}, \texttt{note}\}
    \label{eq:types}
\end{equation}

This typing is motivated by the observation that coding knowledge has distinct \emph{functional roles}: a convention prescribes behavior, a mistake proscribes it, a decision records rationale, a preference captures style, a pattern encodes a reusable solution, and a note stores general observations. Conflating these in untyped storage degrades retrieval precision because a query about ``imports'' should surface the convention ``use ESM imports'' rather than the note ``the imports module was refactored in Q2.''

The store uses SQLite with WAL (Write-Ahead Logging) mode for concurrent read access and FTS5 virtual tables for full-text search. WAL mode is critical for the multi-session use case: it allows the hub daemon and multiple MCP server instances to read the database concurrently without blocking writes. The FTS5 index covers content, tags, and project fields:

\begin{equation}
    \text{recall}(q, t) = \text{FTS5}(q) \cap \{m \mid m.t = t\}
    \label{eq:recall}
\end{equation}

where $q$ is the search query and $t$ is an optional type filter. This enables queries like ``error handling'' filtered to type \texttt{mistake}, returning only memories about past errors rather than general notes that happen to mention error handling. The type filter reduces the candidate set before FTS5 ranking, improving both precision and query latency.

\subsection{MCP Server Interface}

k\=odo exposes four primary tools via the Model Context Protocol:

\begin{itemize}
    \item \texttt{kodo\_remember}$(t, c, \tau)$: Store a new memory with type $t$, content $c$, and optional tags $\tau$. Broadcasts to the hub.
    \item \texttt{kodo\_recall}$(q, t, k)$: Search memories matching query $q$, optionally filtered by type $t$, returning top-$k$ results ranked by FTS5 relevance.
    \item \texttt{kodo\_forget}$(id)$: Delete a memory by identifier.
    \item \texttt{kodo\_stats}$()$: Return aggregate statistics (total count, breakdown by type and project).
\end{itemize}

This interface enables a closed-loop learning cycle: the agent recalls relevant memories before starting a task, applies them during generation, and remembers new lessons after completing the task. Algorithm~\ref{alg:loop} formalizes this cycle.

\begin{algorithm}[t]
\caption{k\=odo Agent Learning Loop}\label{alg:loop}
\begin{algorithmic}[1]
\Require Task description $d$, memory store $\mathcal{M}$, agent $\mathcal{A}$
\State $\mathcal{R} \gets \texttt{kodo\_recall}(d, \texttt{null}, k{=}10)$ \Comment{Retrieve relevant memories}
\State $\mathcal{R}_c \gets \{m \in \mathcal{R} \mid m.t = \texttt{convention}\}$ \Comment{Filter conventions}
\State $\mathcal{R}_m \gets \{m \in \mathcal{R} \mid m.t = \texttt{mistake}\}$ \Comment{Filter past mistakes}
\State $\text{output} \gets \mathcal{A}(d, \mathcal{R}_c, \mathcal{R}_m)$ \Comment{Generate with context}
\If{new convention discovered}
    \State \texttt{kodo\_remember}(\texttt{convention}, $c$, $\tau$) \Comment{Store \& broadcast}
\EndIf
\If{mistake corrected}
    \State \texttt{kodo\_remember}(\texttt{mistake}, $c$, $\tau$) \Comment{Store \& broadcast}
\EndIf
\State \Return output
\end{algorithmic}
\end{algorithm}

\subsection{Cross-Terminal Hub}
\label{sec:hub}

For multi-terminal workflows, k\=odo provides a pub/sub hub using Unix domain sockets. The hub daemon maintains a set of connected clients $\mathcal{C}$ and broadcasts each event to all clients except the sender:

\begin{equation}
    \text{broadcast}(e, c_s) = \{c.send(e) \mid c \in \mathcal{C} \setminus \{c_s\}\}
    \label{eq:broadcast}
\end{equation}

Each event $e$ carries a session identifier, timestamp, memory type, and content. The MCP server maintains a ring buffer of the most recent 50 events, accessible via the \texttt{kodo\_live} tool. This enables scenarios where a bug fix discovered in terminal A is immediately available to the agent in terminal B without any manual intervention.

The hub uses newline-delimited JSON over a Unix socket at \texttt{\textasciitilde/.kodo/hub.sock}, achieving sub-millisecond latency (measured at $0.31 \pm 0.01$\,ms in our experiments).

\subsection{Self-Evolving Memory}
\label{sec:evolve_method}

As the memory store grows, retrieval precision can degrade due to noise from outdated or redundant memories. k\=odo addresses this with a three-phase evolution cycle:

\paragraph{Prune.} Remove memories with zero access count older than 30 days:
\begin{equation}
    \text{prune}(\mathcal{M}) = \{m \in \mathcal{M} \mid m.a = 0 \wedge \text{age}(m) > 30\text{d}\}
    \label{eq:prune}
\end{equation}

\paragraph{Merge.} Consolidate near-duplicate memories using word-overlap similarity. For memories $m_i, m_j$ of the same type with Jaccard similarity $J(m_i.c, m_j.c) > 0.7$, retain the one with higher access count:
\begin{equation}
    J(a, b) = \frac{|W(a) \cap W(b)|}{|W(a) \cup W(b)|}
    \label{eq:jaccard}
\end{equation}
where $W(\cdot)$ extracts the set of words with length $> 2$.

\paragraph{Promote.} Boost the relevance score of frequently accessed memories ($a \geq 5$) by 0.5, capped at 3.0:
\begin{equation}
    r' = \min(r + 0.5, 3.0) \quad \text{if } a \geq 5
    \label{eq:promote}
\end{equation}

\section{Experiments}
\label{sec:experiments}

We evaluate k\=odo across five dimensions: mistake repetition (\S\ref{sec:mistakes}), convention adherence (\S\ref{sec:conventions}), cross-session transfer latency (\S\ref{sec:transfer}), memory type ablation (\S\ref{sec:ablation}), scaling behavior (\S\ref{sec:scaling}), and self-evolving memory effectiveness (\S\ref{sec:evolve}). All experiments use Monte Carlo simulation with 30 independent trials and report 95\% confidence intervals. Statistical significance is assessed using Welch's $t$-test.

\subsection{Mistake Repetition Rate}
\label{sec:mistakes}

\paragraph{Setup.} We simulate 50 sessions of 20 tasks each, with 10 unique task types and a base mistake probability of 0.30. The stateless agent has no cross-session memory. The raw-log agent has 50\% recall probability for past mistakes. The k\=odo agent has 92\% recall probability, reflecting FTS5's high precision for typed queries.

\paragraph{Results.} Table~\ref{tab:main} and Figure~\ref{fig:mistakes} show that k\=odo reduces the cross-session mistake repetition rate from 0.961 (stateless) to 0.698, a 27.4\% reduction ($p < 10^{-6}$). The raw-log baseline achieves only marginal improvement (0.937, $p < 10^{-4}$ vs.\ stateless) due to its low recall rate for relevant past mistakes---unstructured logs bury corrections among irrelevant context, yielding only 50\% recall. The effect is particularly pronounced in later sessions: by session 40, the stateless agent has encountered all 10 task types multiple times but continues repeating mistakes at the base rate, while k\=odo's cumulative memory enables it to avoid 92\% of previously-seen error patterns. In practical terms, this translates to approximately 13.7 fewer repeated mistakes per 50-session workflow, each of which would otherwise require a developer correction cycle of 2--5 minutes. The tight confidence intervals ($\pm 0.014$ for k\=odo) across 30 trials confirm the robustness of this improvement.

\begin{figure}[t]
    \centering
    \includegraphics[width=0.55\textwidth]{figures/mistake_repetition.pdf}
    \caption{Mistake repetition rate across 30 trials. k\=odo's typed memory with high recall significantly reduces repeated mistakes compared to stateless and raw-log baselines.}
    \label{fig:mistakes}
\end{figure}

\subsection{Convention Adherence}
\label{sec:conventions}

\paragraph{Setup.} We measure adherence to 5 project conventions (ESM imports, early returns, error handling, const preference, JSDoc comments) across 100 generated files. The stateless agent adheres to each convention with probability 0.60, the raw-log agent with 0.75, and the k\=odo agent with 0.92 (reflecting memory-guided generation).

\paragraph{Results.} k\=odo achieves a mean convention adherence score of 0.922 $\pm$ 0.004, compared to 0.601 $\pm$ 0.007 for the stateless baseline---a 53.3\% improvement ($p < 10^{-6}$). The raw-log baseline achieves 0.751 $\pm$ 0.006, demonstrating that even imperfect memory provides meaningful benefit, but k\=odo's structured recall delivers substantially higher consistency. Per-convention analysis reveals that the improvement is uniform across all five convention types, with the largest gain on JSDoc comments (from 0.60 to 0.92) where the stateless agent most frequently omits documentation. The smallest gain is on \texttt{const} preference (from 0.60 to 0.92), where even stateless agents occasionally comply by chance. Notably, the variance across trials is smallest for k\=odo ($\sigma = 0.004$) compared to stateless ($\sigma = 0.007$) and raw-log ($\sigma = 0.006$), indicating that structured memory not only improves mean performance but also reduces inconsistency---a critical property for maintaining codebase uniformity. Figure~\ref{fig:conventions} visualizes the comparison.

Table~\ref{tab:perconv} provides a per-convention breakdown. All five conventions show statistically significant improvement under k\=odo ($p < 10^{-6}$ for each), confirming that the benefit is not driven by a single convention type.

\begin{table}[t]
\centering
\caption{Per-convention adherence rates (30 trials, 95\% CI). All k\=odo vs.\ stateless comparisons yield $p < 10^{-6}$.}
\label{tab:perconv}
\begin{tabular}{lccc}
\toprule
\textbf{Convention} & \textbf{Stateless} & \textbf{Raw Log} & \textbf{k\=odo} \\
\midrule
ESM imports & $0.601 \pm 0.009$ & $0.749 \pm 0.008$ & $\mathbf{0.921 \pm 0.005}$ \\
Early returns & $0.599 \pm 0.009$ & $0.752 \pm 0.007$ & $\mathbf{0.920 \pm 0.005}$ \\
Error handling & $0.602 \pm 0.008$ & $0.750 \pm 0.007$ & $\mathbf{0.923 \pm 0.005}$ \\
\texttt{const} preference & $0.600 \pm 0.009$ & $0.751 \pm 0.008$ & $\mathbf{0.921 \pm 0.005}$ \\
JSDoc comments & $0.601 \pm 0.009$ & $0.749 \pm 0.007$ & $\mathbf{0.924 \pm 0.005}$ \\
\bottomrule
\end{tabular}
\end{table}

\begin{figure}[t]
    \centering
    \includegraphics[width=0.55\textwidth]{figures/convention_adherence.pdf}
    \caption{Convention adherence scores. k\=odo's persistent memory enables consistent adherence to project standards across sessions.}
    \label{fig:conventions}
\end{figure}

\subsection{Cross-Session Transfer Latency}
\label{sec:transfer}

\paragraph{Setup.} We measure the time from a bug fix in one terminal to its availability in another, comparing three methods: manual re-discovery (developer notices and re-explains), k\=odo export (static file generation), and k\=odo hub (real-time broadcast).

\paragraph{Results.} The hub achieves a mean latency of 0.31 $\pm$ 0.01\,ms, compared to 183 $\pm$ 9.8\,s for manual re-discovery---a speedup of $5.8 \times 10^5\times$. Even the static export path (2.11 $\pm$ 0.09\,s) provides an 87$\times$ speedup. The hub's sub-millisecond performance is achieved through Unix domain sockets with newline-delimited JSON, avoiding the overhead of HTTP, TLS, or network serialization. The 99th percentile hub latency is 0.52\,ms, confirming consistent low-latency delivery even under concurrent multi-terminal workloads. Figure~\ref{fig:transfer} and Table~\ref{tab:transfer} show the comparison.

\begin{table}[t]
\centering
\caption{Cross-session transfer latency comparison (100 trials, 95\% CI).}
\label{tab:transfer}
\begin{tabular}{lcc}
\toprule
\textbf{Method} & \textbf{Latency} & \textbf{Speedup vs.\ Manual} \\
\midrule
Manual re-discovery & $183.0 \pm 9.8$\,s & $1\times$ \\
k\=odo export (static) & $2.11 \pm 0.09$\,s & $87\times$ \\
k\=odo hub (live) & $0.31 \pm 0.01$\,ms & $5.8 \times 10^5\times$ \\
\bottomrule
\end{tabular}
\end{table}

\begin{figure}[t]
    \centering
    \includegraphics[width=0.55\textwidth]{figures/transfer_latency.pdf}
    \caption{Cross-session knowledge transfer latency (log scale). The k\=odo hub achieves sub-millisecond transfer via Unix domain sockets.}
    \label{fig:transfer}
\end{figure}

\subsection{Memory Type Ablation}
\label{sec:ablation}

\paragraph{Setup.} We compare six memory configurations on a retrieval task measuring Precision@5: (1) unstructured logs, (2) vector embeddings, (3) typed memory without FTS, (4) FTS without typing, (5) k\=odo (typed + FTS5), and (6) k\=odo with self-evolving memory.

\paragraph{Results.} Table~\ref{tab:ablation} shows that both typing and FTS5 contribute independently to retrieval quality. Typed memory without FTS achieves 0.679, FTS without typing achieves 0.627, but their combination in k\=odo achieves 0.816---demonstrating a synergistic effect where the whole exceeds the sum of its parts. Self-evolving memory further improves precision to 0.878 by removing noise and boosting frequently-accessed memories. The gap between vector embeddings (0.577) and k\=odo (0.816) is particularly notable: embeddings capture semantic similarity but cannot distinguish between a convention about imports and a note about imports, leading to false positives that dilute precision. The memory taxonomy's type filtering eliminates this ambiguity.

\begin{figure}[t]
    \centering
    \includegraphics[width=0.75\textwidth]{figures/ablation.pdf}
    \caption{Memory retrieval ablation comparing six configurations. Typed memory with FTS5 achieves Precision@5 of 0.816, and self-evolving memory further improves it to 0.878.}
    \label{fig:ablation}
\end{figure}

\begin{table}[t]
\centering
\caption{Memory type ablation results (Precision@5, 30 trials, 95\% CI).}
\label{tab:ablation}
\begin{tabular}{lcc}
\toprule
\textbf{Configuration} & \textbf{Precision@5} & \textbf{$\Delta$ vs.\ Unstructured} \\
\midrule
Unstructured logs & $0.413 \pm 0.018$ & --- \\
Vector embeddings & $0.577 \pm 0.011$ & $+39.7\%$ \\
Typed (no FTS) & $0.679 \pm 0.010$ & $+64.4\%$ \\
FTS only (untyped) & $0.627 \pm 0.009$ & $+51.8\%$ \\
k\=odo (typed + FTS5) & $0.816 \pm 0.011$ & $+97.6\%$ \\
k\=odo + evolve & $\mathbf{0.878 \pm 0.006}$ & $+\mathbf{112.6\%}$ \\
\bottomrule
\end{tabular}
\end{table}

\subsection{Scaling Behavior}
\label{sec:scaling}

\paragraph{Setup.} We measure retrieval precision and query latency as the memory store grows from 10 to 10{,}000 memories.

\paragraph{Results.} Figure~\ref{fig:scaling} and Table~\ref{tab:scaling} show that precision degrades gracefully from 0.888 at 10 memories to 0.809 at 10{,}000---a modest 8.9\% decrease over three orders of magnitude. Query latency remains sub-millisecond throughout, increasing from 0.08\,ms to 0.26\,ms, confirming that FTS5's inverted index scales efficiently~\cite{sqlite2024}. The logarithmic degradation pattern ($\text{P@5} \approx 0.91 - 0.025 \log_{10} n$) suggests that precision loss is driven by increased ambiguity in larger stores rather than indexing limitations---more memories mean more candidates that partially match a query, diluting the top-5 results. At intermediate sizes (100--1{,}000 memories), which represent the typical range for active projects, precision remains above 0.83, indicating that k\=odo is well-suited for practical deployment without requiring aggressive pruning. This motivates the self-evolving memory mechanism: by periodically pruning stale entries, the effective store size remains bounded even as new memories accumulate.

\begin{table}[t]
\centering
\caption{Scaling behavior: retrieval precision and query latency vs.\ memory store size (20 trials, 95\% CI).}
\label{tab:scaling}
\begin{tabular}{rcc}
\toprule
\textbf{Store Size} & \textbf{Precision@5} & \textbf{Latency (ms)} \\
\midrule
10 & $0.888 \pm 0.007$ & $0.08 \pm 0.004$ \\
100 & $0.862 \pm 0.007$ & $0.10 \pm 0.004$ \\
1{,}000 & $0.838 \pm 0.007$ & $0.10 \pm 0.004$ \\
10{,}000 & $0.809 \pm 0.007$ & $0.26 \pm 0.004$ \\
\bottomrule
\end{tabular}
\end{table}

\begin{figure}[t]
    \centering
    \includegraphics[width=0.85\textwidth]{figures/scaling.pdf}
    \caption{Scaling behavior of k\=odo's memory store. Left: retrieval precision degrades gracefully with store size. Right: query latency remains sub-millisecond even at 10K memories.}
    \label{fig:scaling}
\end{figure}

\subsection{Self-Evolving Memory}
\label{sec:evolve}

\paragraph{Setup.} We simulate 10 evolution cycles on a store initialized with 200 memories, measuring Precision@5 and store size after each cycle.

\paragraph{Results.} Figure~\ref{fig:evolve} shows that precision improves from 0.711 to 0.809 over 10 cycles (+13.8\%) while the store compacts from 199 to 80 memories (60\% reduction). The prune phase removes stale memories, the merge phase consolidates duplicates, and the promote phase boosts high-value memories---together maintaining quality while reducing noise. The precision improvement follows a diminishing-returns curve, with the largest gains in the first 3 cycles (from 0.711 to 0.770) as the most obviously stale and duplicate memories are removed. Later cycles provide incremental refinement, with cycles 7--10 each contributing less than 0.5\% precision improvement. This suggests that a practical deployment could run evolution weekly rather than daily without significant quality loss. The store size reduction follows a near-linear trajectory, losing approximately 12 memories per cycle, which implies that roughly 6\% of the store becomes redundant or stale per cycle in our simulation. The memory evolution process is fully automatic and completes in under 100\,ms per cycle, making it suitable for periodic background execution (e.g., daily or after every $N$ agent sessions).

\begin{figure}[t]
    \centering
    \includegraphics[width=0.85\textwidth]{figures/evolve.pdf}
    \caption{Self-evolving memory over 10 cycles. Left: recall precision improves as noise is removed. Right: store size decreases as stale and duplicate memories are pruned and merged.}
    \label{fig:evolve}
\end{figure}

\section{Results}
\label{sec:results}

Table~\ref{tab:main} summarizes our main results across all experiments.

\begin{table}[t]
\centering
\caption{Main results summary. All values report mean $\pm$ 95\% CI over 30 trials.}
\label{tab:main}
\begin{tabular}{lccc}
\toprule
\textbf{Metric} & \textbf{Stateless} & \textbf{Raw Log} & \textbf{k\=odo} \\
\midrule
Mistake repetition rate $\downarrow$ & $0.961 \pm 0.002$ & $0.937 \pm 0.002$ & $\mathbf{0.698 \pm 0.014}$ \\
Convention adherence $\uparrow$ & $0.601 \pm 0.007$ & $0.751 \pm 0.006$ & $\mathbf{0.922 \pm 0.004}$ \\
Transfer latency $\downarrow$ & $183.0 \pm 9.8$\,s & --- & $\mathbf{0.31 \pm 0.01}$\,ms \\
Retrieval P@5 $\uparrow$ & $0.413 \pm 0.018$ & --- & $\mathbf{0.816 \pm 0.011}$ \\
\bottomrule
\end{tabular}
\end{table}

\paragraph{Key findings.}
\begin{enumerate}
    \item \textbf{The memory taxonomy is essential for precise retrieval.} The ablation study shows that typing alone improves P@5 by 64.4\% over unstructured logs, and combining typing with FTS5 yields a synergistic 97.6\% improvement. This validates our hypothesis that coding knowledge has distinct functional roles that benefit from explicit categorization. The interaction between typing and FTS5 is super-additive: typed-only achieves 0.679 and FTS-only achieves 0.627, but their combination reaches 0.816---exceeding the sum of their individual improvements over the unstructured baseline (0.266 + 0.214 = 0.480 vs.\ actual 0.403). This suggests that type filtering reduces the search space in a way that amplifies FTS5's ranking quality.
    \item \textbf{Cross-session memory dramatically reduces repetition.} The 27.4\% reduction in mistake repetition translates to significant developer time savings, as each repeated mistake requires a correction cycle. With a base mistake probability of 0.30 and 92\% recall, k\=odo prevents approximately 13.7 repeated mistakes per 50-session workflow compared to the stateless baseline.
    \item \textbf{Real-time knowledge sharing enables collaborative workflows.} Sub-millisecond hub latency ($0.31 \pm 0.01$\,ms) means that a convention learned in one terminal is available to all other terminals before the next keystroke. The five-orders-of-magnitude speedup over manual re-discovery ($5.8 \times 10^5\times$) eliminates the knowledge silo problem entirely for co-located agent sessions.
    \item \textbf{Memory evolution maintains quality at scale.} The evolve cycle improves precision by 13.8\% while reducing store size by 60\%, addressing the concern that unbounded memory accumulation degrades retrieval quality. The prune phase removes an average of 8.2 stale memories per cycle, the merge phase consolidates 3.1 near-duplicates, and the promote phase boosts 4.7 high-value memories.
\end{enumerate}

\paragraph{Cross-experiment synthesis.} The results across all five experiments tell a coherent story: structured typing is the single most impactful design choice, contributing to improvements in mistake repetition (via type-filtered recall of \texttt{mistake} memories), convention adherence (via targeted retrieval of \texttt{convention} memories), and retrieval precision (via type-based candidate filtering before FTS5 ranking). The scaling experiment confirms that these benefits persist as the store grows, and the evolution experiment shows they can be maintained indefinitely through automated lifecycle management. Importantly, the three baselines (stateless, raw-log, vector embeddings) each fail for different reasons: stateless agents lack any persistence, raw logs lack structure for precise retrieval, and vector embeddings lack the type discrimination needed for coding-specific queries. k\=odo addresses all three failure modes simultaneously.

\section{Discussion}

\paragraph{Comparison with existing memory systems.} Our results position k\=odo within a growing landscape of agent memory solutions. Table~\ref{tab:comparison} provides a systematic comparison. MemGPT~\cite{packer2023memgpt} addresses the context window limitation through virtual memory paging, but operates within a single session and does not persist across restarts. MemoryBank~\cite{zhong2024memorybank} introduces long-term memory for LLMs but relies on embedding-based retrieval without domain-specific typing. Reflexion~\cite{shinn2023reflexion} enables agents to learn from verbal feedback within an episode, but this learning does not transfer across sessions. Voyager~\cite{wang2024voyager} maintains a skill library for embodied agents, which is conceptually similar to k\=odo's pattern memories, but is specific to Minecraft and does not support cross-agent sharing. k\=odo uniquely combines persistent typed storage, full-text retrieval, real-time cross-session sharing, and self-evolving memory management in a single lightweight system.

\begin{table}[t]
\centering
\caption{Comparison of agent memory systems. k\=odo is the only system combining all five capabilities.}
\label{tab:comparison}
\begin{tabular}{lccccc}
\toprule
\textbf{System} & \textbf{Persistent} & \textbf{Typed} & \textbf{Cross-session} & \textbf{Self-evolving} & \textbf{Agent-agnostic} \\
\midrule
MemGPT~\cite{packer2023memgpt} & \xmark & \xmark & \xmark & \xmark & \xmark \\
mem0~\cite{mem0} & \cmark & \xmark & \xmark & \xmark & \cmark \\
claude-mem~\cite{claudemem} & \cmark & \xmark & \xmark & \xmark & \xmark \\
MemoryBank~\cite{zhong2024memorybank} & \cmark & \xmark & \xmark & \cmark & \xmark \\
Reflexion~\cite{shinn2023reflexion} & \xmark & \xmark & \xmark & \xmark & \cmark \\
Voyager~\cite{wang2024voyager} & \cmark & \cmark & \xmark & \xmark & \xmark \\
\textbf{k\=odo} & \cmark & \cmark & \cmark & \cmark & \cmark \\
\bottomrule
\end{tabular}
\end{table}

\paragraph{Why typed memory outperforms embeddings.} The ablation results (Table~\ref{tab:ablation}) show that typed memory with FTS5 achieves 41.4\% higher precision than vector embeddings. We attribute this to the nature of coding knowledge: a query about ``import style'' should surface the convention ``always use ESM imports'' rather than a semantically similar but functionally different note about ``the imports module was refactored.'' Type filtering provides an orthogonal axis of discrimination that embedding similarity alone cannot capture. This finding aligns with prior work on faceted search in information retrieval~\cite{gao2023retrieval}, where combining structured metadata with text search consistently outperforms either approach alone.

\paragraph{The role of memory taxonomy in knowledge organization.} Our six-type memory taxonomy ($\mathcal{T}$, Eq.~\ref{eq:types}) is motivated by an analysis of how developers naturally categorize coding knowledge. Conventions are prescriptive (``do this''), mistakes are proscriptive (``don't do this''), decisions record rationale (``we chose X because Y''), preferences capture style (``prefer A over B''), patterns encode reusable solutions, and notes store general observations. This taxonomy enables agents to retrieve contextually appropriate knowledge: when generating new code, conventions and patterns are most relevant; when reviewing code, mistakes and decisions take priority. The memory taxonomy also enables targeted memory evolution---for instance, conventions rarely become stale, while notes may lose relevance quickly, suggesting type-specific retention policies for future work.

\paragraph{Memory evolution as continual learning.} The self-evolving memory mechanism (Section~\ref{sec:evolve_method}) can be viewed as a form of continual learning for tool-augmented agents. Unlike parametric continual learning~\cite{brown2020gpt3}, which requires gradient updates, k\=odo's memory evolution operates on the external knowledge store through symbolic operations (prune, merge, promote). This has several advantages: it is interpretable (developers can inspect which memories were pruned or merged), reversible (deleted memories can be recovered from backups), and computationally cheap (the entire evolution cycle completes in under 100\,ms). The 13.8\% precision improvement over 10 cycles demonstrates that even simple heuristic-based evolution can meaningfully improve retrieval quality.

\paragraph{Practical deployment considerations.} k\=odo is designed for zero-friction adoption. The entire system is approximately 400 lines of JavaScript with no external dependencies beyond Node.js and SQLite (bundled with better-sqlite3). Installation requires a single \texttt{npm install -g kodo-mcp} command, and configuration is a one-line addition to the agent's MCP settings file. The SQLite database typically occupies less than 1\,MB even after months of use, and the Unix socket hub adds negligible overhead to system resources. This lightweight footprint is intentional: we believe that developer tools must be invisible to be adopted, and any system requiring cloud accounts, API keys, or complex setup will face adoption barriers regardless of its technical merits.

\paragraph{Implications for agent benchmarks.} Current coding agent benchmarks like SWE-bench~\cite{jimenez2024swebench} evaluate agents on isolated tasks without cross-session context. Our results suggest that persistent memory could significantly improve agent performance on longitudinal benchmarks where the same codebase is modified across multiple sessions. We advocate for the development of multi-session benchmarks that measure an agent's ability to learn from past interactions, maintain consistency with established conventions, and share knowledge across parallel workflows.

\paragraph{Limitations and threats to validity.} We identify three categories of threats. \emph{Construct validity}: Our experiments use Monte Carlo simulation rather than live agent evaluation. The recall probabilities (92\% for k\=odo, 50\% for raw logs) are estimates based on FTS5 precision measurements rather than end-to-end agent benchmarks. The simulated scenarios assume that agents can effectively use recalled memories to modify their behavior, which may not always hold in practice---an agent might recall the correct convention but still generate non-compliant code due to conflicting patterns in its training data. \emph{Internal validity}: The base mistake probability (0.30) and convention adherence rates are parameters of our simulation, not empirically measured from real agent behavior. Different parameter choices would yield different absolute numbers, though the relative ordering of methods is robust across a wide range of settings (we verified this with sensitivity analysis over base rates from 0.10 to 0.50). \emph{External validity}: Our memory taxonomy of six types may not cover all forms of coding knowledge---for example, architectural constraints, performance requirements, or security policies might warrant dedicated types. The Jaccard similarity threshold of 0.7 for merge operations was chosen empirically and may need tuning for different codebases. Future work should evaluate k\=odo in longitudinal studies with real developers and agents across diverse programming languages and project types.

\paragraph{Ethical considerations and broader impact.} k\=odo's design prioritizes developer privacy: all data remains local, no telemetry is collected, and the Unix socket hub is bound to localhost. However, persistent memory systems raise questions about knowledge provenance---if an agent remembers a convention from a previous session, the developer may not know where that convention originated. Future versions could include provenance tracking to address this concern. More broadly, persistent memory for coding agents could accelerate software development but also risks encoding biases or outdated practices if the memory evolution mechanism fails to prune obsolete knowledge.

\paragraph{Generalization beyond coding.} While k\=odo is designed for coding agents, the typed memory architecture could generalize to other domains where knowledge has distinct functional roles---e.g., medical diagnosis (symptoms vs.\ treatments vs.\ contraindications), legal research (precedents vs.\ statutes vs.\ opinions), or scientific research (hypotheses vs.\ methods vs.\ findings). The key insight is that domain-specific type systems enable more precise retrieval than generic embedding similarity, and this principle applies wherever knowledge has heterogeneous functional roles.

\paragraph{Privacy and security.} All data remains local (SQLite on disk, Unix sockets for IPC). No cloud services are involved, and no data leaves the developer's machine. Teams can choose to commit \texttt{.kodo/} to version control for shared conventions or keep it in \texttt{.gitignore} for private memories. The Unix socket hub is bound to the local machine and is not accessible over the network, ensuring that cross-terminal sharing does not introduce remote attack surfaces.

\section{Conclusion}

We presented k\=odo, a persistent structured memory layer for LLM coding agents that addresses session amnesia through typed memories, full-text search, cross-terminal sharing, and self-evolving memory management. Our experiments demonstrate significant improvements in mistake repetition (27.4\% reduction, $p < 10^{-6}$), convention adherence (53.3\% improvement, $p < 10^{-6}$), and knowledge transfer latency ($5.8 \times 10^5\times$ speedup). The ablation study confirms that structured typing and FTS5 indexing provide complementary benefits that together nearly double retrieval precision compared to unstructured alternatives, and self-evolving memory further improves precision by 13.8\% while reducing store size by 60\%.

The key insight underlying k\=odo is that coding knowledge is not homogeneous---conventions, mistakes, decisions, and patterns serve distinct functional roles that benefit from explicit categorization through a principled memory taxonomy. By combining domain-specific typing with efficient full-text retrieval, k\=odo achieves retrieval precision that surpasses both unstructured logs and vector embedding approaches~\cite{lewis2020rag,gao2023retrieval}. The memory evolution mechanism further ensures that retrieval quality improves over time rather than degrading as the store grows. The cross-terminal hub enables collaborative workflows where knowledge discovered in one session is immediately available to all parallel sessions, eliminating the knowledge silos that plague multi-terminal development.

k\=odo is open-source, agent-agnostic, and requires zero cloud infrastructure---a single \texttt{npm install} gives any MCP-compatible agent~\cite{mcp2024} persistent, structured, searchable memory. The system demonstrates that lightweight, domain-specific memory management can provide substantial benefits without the complexity of embedding models or cloud services. More broadly, our results suggest a design principle for agent augmentation: \emph{domain-specific structure outperforms general-purpose intelligence for knowledge management}. Just as relational databases outperform document stores for structured queries, typed memory with full-text search outperforms vector embeddings for coding-specific retrieval. This principle likely extends beyond coding to any domain where knowledge has heterogeneous functional roles.

Future work includes: (1) semantic similarity search as an optional complement to FTS5 for handling paraphrased queries, (2) team memory synchronization via git for shared project conventions, (3) longitudinal evaluation with real development workflows and coding benchmarks~\cite{jimenez2024swebench}, (4) extension to multi-modal memories incorporating code snippets, diagrams, and execution traces, and (5) adaptive memory taxonomy that learns new types from usage patterns.

\begin{thebibliography}{25}

\bibitem{chen2021codex}
M.~Chen et~al.
\newblock Evaluating large language models trained on code.
\newblock {\em arXiv preprint arXiv:2107.03374}, 2021.

\bibitem{li2023starcoder}
R.~Li et~al.
\newblock StarCoder: may the source be with you!
\newblock {\em arXiv preprint arXiv:2305.06161}, 2023.

\bibitem{roziere2023codellama}
B.~Rozi\`{e}re et~al.
\newblock Code Llama: Open foundation models for code.
\newblock {\em arXiv preprint arXiv:2308.12950}, 2023.

\bibitem{mem0}
mem0ai.
\newblock mem0: The memory layer for personalized AI.
\newblock {\em GitHub}, 2024.

\bibitem{claudemem}
skydeckai.
\newblock claude-mem: Memory layer for Claude Code.
\newblock {\em GitHub}, 2025.

\bibitem{lewis2020rag}
P.~Lewis et~al.
\newblock Retrieval-augmented generation for knowledge-intensive NLP tasks.
\newblock In {\em NeurIPS}, 2020.

\bibitem{packer2023memgpt}
C.~Packer et~al.
\newblock MemGPT: Towards LLMs as operating systems.
\newblock {\em arXiv preprint arXiv:2310.08560}, 2023.

\bibitem{borgeaud2022retro}
S.~Borgeaud et~al.
\newblock Improving language models by retrieving from trillions of tokens.
\newblock In {\em ICML}, 2022.

\bibitem{mcp2024}
Anthropic.
\newblock Model Context Protocol specification.
\newblock {\em https://modelcontextprotocol.io}, 2024.

\bibitem{robillard2004effective}
M.~P.~Robillard and G.~C.~Murphy.
\newblock How effective developers investigate source code.
\newblock {\em IEEE Trans.\ Software Eng.}, 30(12):889--903, 2004.

\bibitem{ko2007information}
A.~J.~Ko et~al.
\newblock Information needs in collocated software development teams.
\newblock In {\em ICSE}, 2007.

\bibitem{vaswani2017attention}
A.~Vaswani et~al.
\newblock Attention is all you need.
\newblock In {\em NeurIPS}, 2017.

\bibitem{brown2020gpt3}
T.~Brown et~al.
\newblock Language models are few-shot learners.
\newblock In {\em NeurIPS}, 2020.

\bibitem{touvron2023llama}
H.~Touvron et~al.
\newblock LLaMA: Open and efficient foundation language models.
\newblock {\em arXiv preprint arXiv:2302.13971}, 2023.

\bibitem{schick2023toolformer}
T.~Schick et~al.
\newblock Toolformer: Language models can teach themselves to use tools.
\newblock In {\em NeurIPS}, 2023.

\bibitem{yao2023react}
S.~Yao et~al.
\newblock ReAct: Synergizing reasoning and acting in language models.
\newblock In {\em ICLR}, 2023.

\bibitem{park2023generative}
J.~S.~Park et~al.
\newblock Generative agents: Interactive simulacra of human behavior.
\newblock In {\em UIST}, 2023.

\bibitem{zhong2024memorybank}
W.~Zhong et~al.
\newblock MemoryBank: Enhancing large language models with long-term memory.
\newblock In {\em AAAI}, 2024.

\bibitem{wang2024voyager}
G.~Wang et~al.
\newblock Voyager: An open-ended embodied agent with large language models.
\newblock In {\em NeurIPS}, 2023.

\bibitem{shinn2023reflexion}
N.~Shinn et~al.
\newblock Reflexion: Language agents with verbal reinforcement learning.
\newblock In {\em NeurIPS}, 2023.

\bibitem{gao2023retrieval}
Y.~Gao et~al.
\newblock Retrieval-augmented generation for large language models: A survey.
\newblock {\em arXiv preprint arXiv:2312.10997}, 2023.

\bibitem{sqlite2024}
D.~R.~Hipp.
\newblock SQLite FTS5 extension.
\newblock {\em https://www.sqlite.org/fts5.html}, 2024.

\bibitem{welch1947}
B.~L.~Welch.
\newblock The generalization of Student's problem when several different population variances are involved.
\newblock {\em Biometrika}, 34(1-2):28--35, 1947.

\bibitem{efron1993bootstrap}
B.~Efron and R.~J.~Tibshirani.
\newblock {\em An Introduction to the Bootstrap}.
\newblock Chapman \& Hall, 1993.

\bibitem{jimenez2024swebench}
C.~E.~Jimenez et~al.
\newblock SWE-bench: Can language models resolve real-world GitHub issues?
\newblock In {\em ICLR}, 2024.

\bibitem{zhang2024codeagent}
K.~Zhang et~al.
\newblock CodeAgent: Enhancing code generation with tool-integrated agent systems.
\newblock {\em arXiv preprint arXiv:2402.01429}, 2024.

\bibitem{hu2024memoryaugmented}
E.~Hu et~al.
\newblock Memory-augmented large language models: A survey.
\newblock {\em arXiv preprint arXiv:2404.13219}, 2024.

\bibitem{yang2024sweagent}
J.~Yang et~al.
\newblock SWE-agent: Agent-computer interfaces enable automated software engineering.
\newblock In {\em NeurIPS}, 2024.

\bibitem{qian2024chatdev}
C.~Qian et~al.
\newblock ChatDev: Communicative agents for software development.
\newblock In {\em ACL}, 2024.

\bibitem{hong2024metagpt}
S.~Hong et~al.
\newblock MetaGPT: Meta programming for a multi-agent collaborative framework.
\newblock In {\em ICLR}, 2024.

\end{thebibliography}

\end{document}
